\documentclass[11pt]{article}
\usepackage[margin=1in]{geometry}
\usepackage{listings}
\usepackage{xcolor}
\usepackage{hyperref}

\lstset{
  basicstyle=\small\ttfamily,
  breaklines=true,
  frame=single,
  backgroundcolor=\color{gray!10}
}

\title{Supplementary Material: Monte Carlo Simulation Code\\
\large Eating the Fire: Biomimetic Boundary Layers and the\\
Case for Radiation Harvesting at Nuclear Interfaces}
\author{Leslie Yarbrough}
\date{April 2026}

\begin{document}
\maketitle

\section*{Simulation Overview}
The thin-foil gamma harvesting simulator models photon transport through
layered absorber stacks at 2~MeV. Three interaction mechanisms are modeled:
Compton scattering, pair production, and photoelectric absorption.
Interaction type is selected probabilistically from linear attenuation
coefficients for tungsten. Charge collection is determined by whether the
resulting secondary electron can reach a layer boundary within its
approximate range of 10~$\mu$m in metal.

The simulator is available at:
\url{https://github.com/lyarbrough-lab/Vela3}

\section*{Key Results}
\begin{center}
\begin{tabular}{rrrr}
\hline
Layer thickness ($\mu$m) & No.\ layers & Collection (\%) & vs.\ bulk \\
\hline
26000 (bulk) & 1    & $\approx 0$  & 1$\times$ \\
5000         & 5    & 0.017        & 17$\times$ \\
1000         & 26   & 0.030        & 30$\times$ \\
500          & 52   & 0.101        & 101$\times$ \\
100          & 260  & 0.372        & 372$\times$ \\
50           & 520  & 0.933        & 933$\times$ \\
10           & 2600 & 3.391        & $>3000\times$ \\
5            & 5200 & 3.884        & $>3000\times$ \\
\hline
\end{tabular}
\end{center}

Total stack thickness held constant at 2.6~cm across all runs.

\end{document}
